\section*{第 \ref{ch:b2:laser} 章——激光：受激辐射与高斯光束}

\begin{solution}{exo:b2:laser:1}
$\nu = \qty{4.74e14}{Hz}$; $h\nu = \qty{3.14e-19}{J} = \qty{1.96}{eV}$;
每秒 $\qty{3.2e15}{}$ 个光子；$c/2L = \qty{500}{MHz}$；三
个模式。
\end{solution}

\begin{solution}{exo:b2:laser:2}
$\eu^{-80} \approx 2 \times 10^{-35}$；仅当 $T \to \infty$ 时等号成立；
$N_2 > N_1$ 要求 $\eu^{-h\nu/k_BT} > 1$，即形式上 $T < 0$。
当 \qty{24}{GHz} 时，$h\nu = \qty{1e-4}{eV}$：$N_2/N_1 = 0.996$。
\end{solution}

\begin{solution}{exo:b2:laser:3}
$z_{\text{R}} = \qty{1.24}{m}$; $\theta = \qty{4.0e-4}{rad}$; $w(\qty{100}{m}) =
\qty{40}{mm}$；在月球处为 \qty{150}{km}。经 $\times 20$ 后：$w_0 = \qty{10}{mm}$，
$\theta = \qty{2e-5}{rad}$，$z_{\text{R}} = \qty{500}{m}$，在
\qty{100}{m} 处为 \qty{10.2}{mm}，在月球处为 \qty{7.7}{km}。
\end{solution}

\begin{solution}{exo:b2:laser:4}
$g_{\text{阈}} = (0.005 + 0.0101)/0.2 = \qty{0.076}{m^{-1}} < 0.2$：激光器起振。
取 $R_2 = 0.90$ 时：$(0.005 + 0.053)/0.2 = \qty{0.29}{m^{-1}}$：
不起振。
\end{solution}

\begin{solution}{exo:b2:laser:5}
(a) 见\cref{prop:b2:laser:einstein}。(b) \qty{600}{nm}：$h\nu/k_BT
= 80$，比值 $5 \times 10^{34}$；\qty{1}{cm}：$4.8 \times 10^{-3}$——
受激辐射占主导。(c) 在微波频段
自发辐射可以忽略，少量粒子数反转即可
获得增益；所需技术（谐振腔、波导）早已存在。(d) $A \propto
\nu^3$，因此 $\tau \propto \lambda^3$：\qty{10}{\micro s}。
\end{solution}

\begin{solution}{exo:b2:laser:6}
(a) $I_{\text{s}} = h\nu/\sigma\tau = \qty{1.0e5}{W/m^2}$. (b) $R =
g_0/\sigma\tau = \qty{3.3e22}{m^{-3}.s^{-1}}$. (c) $I = I_{\text{s}}(5 - 1) =
\qty{4.2e5}{W/m^2}$; $P_{\text{出}} = 0.01 \times 4.2 \times 10^5 \times
10^{-6} = \qty{4.2}{mW}$. (d) $I = I_{\text{s}}(\sigma R\tau/g_{\text{阈}}
- 1) = (h\nu/g_{\text{阈}})(R - g_{\text{阈}}/\sigma\tau)$：关于 $R$ 呈线性
在 $R_{\text{阈}} = g_{\text{阈}}/\sigma\tau$ 以上。
\end{solution}

\begin{solution}{exo:b2:laser:7}
(a) $c/2L > \qty{1.5}{GHz}$: $L < \qty{10}{cm}$. (b) $\delta\nu = \nu\,\delta L/L
= \qty{1.6}{GHz}$：三个纵模间隔，即整个增益宽度——该模式
扫过增益曲线，并交接给相邻的纵模。(c) 用
加热器将腔长控制在 $\lambda$ 的很小一部分内，以
两纵模的平衡作为误差信号。(d) $c/\delta\nu \approx
\qty{100}{km}$.
\end{solution}

\begin{solution}{exo:b2:laser:8}
(a) $w_0' = \qty{10}{\micro m}$, $z_{\text{R}} = \qty{0.5}{mm}$, $I = 2P/\pi
w_0'^2 = \qty{6.4e6}{W/m^2}$. (b) $w_0' = \qty{3.4}{\micro m}$, $I =
\qty{5.5e7}{W/m^2}$；太阳的像：$10^3 \times (3/0.15)^2 =
\qty{4e5}{W/m^2}$——激光笔差一百倍。(c) 灯泡的
功率分布于 $4\pi$ 立体角内，其在视网膜上的像是扩展的；激光笔的
全部功率集中在几微米之内，位于
最敏感的波长处。(d) $2z_{\text{R}} = \qty{1}{mm}$。
\end{solution}

\begin{solution}{exo:b2:laser:9}
(a) $(N/2)\,h\nu/\tau = \qty{7.6e8}{W/m^3}$。(b) \qty{760}{W}——连续
运行不可行：应采用闪光灯。(c) $10^{22} \times 1.87 \times
10^{-19}/2.3 \times 10^{-4} = \qty{8e6}{W/m^3}$：\qty{8}{W} 每立方
厘米。(d) 功率小一百倍，且为连续工作：
几瓦的二极管光即可满足要求。
\end{solution}

\begin{solution}{exo:b2:laser:10}
(a) $\lambda/\pi w_0$：\qty{0.17}{rad}（\qty{9.5}{\degree}）和 \qty{0.5}{rad}
（\qty{28}{\degree}）。(b) 窄方向的发散更大：
远场椭圆与发射体垂直。(c) $f\theta$：
\qty{0.8}{mm} 和 \qty{2.5}{mm}。(d) 一对变形棱镜或
柱面透镜；或耦合进单模光纤。
\end{solution}

\begin{solution}{exo:b2:laser:11}
(a) 强度每 \qty{2}{ns} 增至 $\times 1.04$ 倍。(b) $N = P\,(2L/c)/h\nu =
6.4 \times 10^8$ 个光子；$\ln(6.4 \times 10^8)/\ln 1.04 \approx 520$ 次
往返，约 \qty{1}{\micro s}。(c) 增益饱和；低于
阈值的模式每次往返损耗大于增益，并停留在
自发辐射水平。(d) 各模式共享同一粒子数反转；最强的
模式使其饱和，使其余模式无增益可用；腔长与增益的漂移
会改变胜出的模式——输出在各模式间跳变。
\end{solution}

\begin{solution}{exo:b2:laser:12}
(a) $2.7 \times 10^{17}$ 个光子；\qty{1}{GW}。(b) $\theta = \qty{3.4e-7}{rad}$，
\qty{130}{m}；视宁度为 \qty{1.8}{km}。(c) 分数为 $0.1/\pi(1800)^2 =
10^{-8}$：$2.6 \times 10^9$ 个光子；返回光的扩散 \qty{1.3e-5}{rad}，半径
为 \qty{5}{km}；望远镜接收到 $(0.5/5000)^2 = 10^{-8}$：约 $25$
个光子，经光学元件后仅剩几个。(d) $c \times \qty{100}{ps}/2 = \qty{1.5}{cm}$；
$\times 1/\sqrt{10^4}$: \qty{0.15}{mm}.
\end{solution}

\begin{solution}{pb:b2:laser:1}
\textbf{1.} \qty{1.96}{eV}, \qty{4.74e14}{Hz}.

\textbf{2.} $\qty{0.05}{eV} \approx 1.5\,k_BT$：热运动提供
微小的亏缺；转移几乎共振，因此效率高。

\textbf{3.} $\eu^{-20.66/0.0345} = \eu^{-600} \approx 10^{-260}$：无；只有
放电向该能级提供粒子。

\textbf{4.} 下能级的排空速度比上能级的
衰变快十倍：它始终几乎为空——即四能级方案。

\textbf{5.} $v^* = \sqrt{2k_BT/m} = \qty{580}{m/s}$; $\nu v^*/c \approx \qty{0.9}{GHz}$;
精确公式给出 \qty{1.5}{GHz}。

\textbf{6.} $g_0 = \qty{0.09}{m^{-1}}$；$\eu^{0.027}$：$2.7\%$ 每程。

\textbf{7.} $g_{\text{阈}} = -\ln(0.999 \times 0.99)/0.6 = \qty{0.018}{m^{-1}}$;
裕量 $5$。

\textbf{8.} \qty{500}{MHz}；三个。

\textbf{9.} $P_{\text{入}} = \qty{1}{mW}/0.01 = \qty{100}{mW}$; \qty{1.3e5}{W/m^2}.

\textbf{10.} 增益饱和至 $g_{\text{阈}}$；多余的
能量转为输出光子（以及自发辐射、热量）。

\textbf{11.} $\qty{1}{mW}/\qty{7.5}{W} = 1.3 \times 10^{-4}$：大部分
电子能量流向氦亚稳态之外，且
光子携带了投入的 \qty{20.66}{eV} 中的 \qty{1.96}{eV}。

\textbf{12.} $2 \times 0.072 \times 0.3 = 0.043$; $N = 0.1 \times 2 \times 10^{-9}/3.14
\times 10^{-19} = 6.4 \times 10^8$；$20.3/0.043 \approx 470$ 次往返；
\qty{0.9}{\micro s}.

\textbf{13.} $\delta\nu = \nu\,\delta L/L = \qty{1.6}{GHz}$：模式穿过
整条增益曲线——除非腔长保持稳定，激光器
就会跳到另一个模式。

\textbf{14.} $p$ 偏振无
反射地穿过布儒斯特表面；$s$ 偏振每个表面损失约 $15\%$，
从不达到阈值：输出线偏振光，且无额外代价。

\textbf{15.} $z_{\text{R}} = \qty{0.45}{m}$; $\theta = \qty{6.7e-4}{rad}$; $w(\qty{10}{m})
= \qty{6.7}{mm}$：全宽 \qty{13}{mm}。

\textbf{16.} $2P/\pi w_0^2 = \qty{7}{kW/m^2}$：相当于七个太阳。

\textbf{17.} $w_0' = \qty{34}{\micro m}$; \qty{5.6e5}{W/m^2}.

\textbf{18.} $w_0' = \qty{11}{\micro m}$；\qty{5e6}{W/m^2}；视网膜能承受
\qty{1}{mW} 持续回避反射所需的 \qty{0.25}{s}——二类：可见，
$\le \qty{1}{mW}$，因为人会眨眼，所以是安全的。

\textbf{19.} \qty{6.7e-5}{rad}; $w_0 = \qty{3}{mm}$, $z_{\text{R}} = \qty{45}{m}$,
$w(\qty{1}{km}) = \qty{67}{mm}$: \qty{13}{cm}.

\textbf{20.} $c/\Delta\nu$: \qty{300}{m}; \qty{30}{cm}.

\textbf{21.} $2.7 \times 10^{17}$; \qty{1}{GW}.

\textbf{22.} \qty{1.8}{km}；仅衍射时为 \qty{130}{m}。

\textbf{23.} $0.1/\pi(1800)^2 = 10^{-8}$；$2.6 \times 10^9$ 个光子。

\textbf{24.} $\lambda/d = \qty{1.3e-5}{rad}$, \qty{5}{km}; $(0.5/5000)^2 = 10^{-8}$:
$25$ 个光子，两个或三个被探测到。

\textbf{25.} \qty{1.5}{cm}；\qty{0.15}{mm}；退行在
数周内测得。
\end{solution}
